\chapter{理论基础与文献综述}
%%====================================================================

\section{开源软件开发与开发者激励}

\subsection{开源软件的特征与发展历程}

开源软件（Open-Source Software, OSS）以源代码公开为基本特征，
允许任何人自由查看、使用、修改和分发\cite{lakhani2003open}。
有别于商业软件的封闭开发模式，开源项目在本质上依托互联网展开分布式协作——
参与者分布于全球各类机构，多数不受雇佣合同约束，
以自愿形式贡献代码、文档和讨论。协调工作则依靠版本控制系统、
邮件列表及代码审查流程等技术与社会机制得以维系。

自由软件运动的理念基础由 Richard Stallman\footnote{Richard Stallman 于1983年发起 GNU 项目，并于1985年创立自由软件基金会（Free Software Foundation, FSF），提出了“自由软件”的概念与 GPL 许可证体系。} 在20世纪80年代奠定；
1991年 Linus Torvalds 发布 Linux 内核，将该理念推向工程实践的前沿。
21世纪以来，互联网基础设施快速普及，加之 Git\footnote{Git 是 Linus Torvalds 于2005年为管理 Linux 内核源代码而开发的分布式版本控制系统，目前已成为软件工程领域的标准协作工具。} 版本控制工具日趋成熟，
全球开发者协作的组织成本大幅下降，以 GitHub 为核心的开源生态由此成形。
如今，Linux 内核支撑着全球绝大多数服务器与移动设备，
Python 已成为数据科学和人工智能研究的通用语言——
这些产品无一不是数以千计志愿开发者多年协作的成果。

从经济学角度审视，开源软件具备公共品的典型特征：
非竞争性与非排他性叠加，潜在的“搭便车”问题由此产生——
理论上人人都倾向于坐享他人贡献，而不愿自行承担成本\cite{lerner2002simple}。
但开源生态的实际状况与上述悲观预测相悖：
大量高质量项目维持了数十年的持续演进，贡献者遍布全球。
这种“反常”现象持续激发学界对开源开发者动机机制的探索\cite{gerosa2021shifting}，
并催生了一套有别于传统经济人假设的理论体系。

开源软件的经济重要性在近十余年间经历了质的跃升。
据估计，企业核心软件产品中超过70\%的代码来源于开源组件，
全球开源软件每年为商业应用贡献的价值已达数千亿美元规模。
然而，支撑这一巨大价值的开发者群体在规模上极度有限，
且整体处于志愿贡献状态——许多关键项目的核心维护工作仅集中在个位数的开发者身上\cite{lu2021core,jamieson2024predicting}。
“巨大经济价值与极少志愿维护者”之间的结构性不对称，
使理解和维持开发者参与动机成为兼具学术价值与现实紧迫性的议题。
近年来，“开源可持续性”这一概念被学界广泛用以描述上述挑战，
研究者呼吁在机构设计和激励机制层面作出系统性应对\cite{alami2024sustainability,wang2025opensource,qi2021opensource}。
本研究正是在该宏观背景下，对开源激励机制中尚未被充分探讨的
信息维度予以回应。

\subsection{开发者参与动机的理论框架}

开源开发者动机研究的核心挑战在于为何大量技能成熟的开发者
愿意将大量时间投入缺乏直接经济回报的项目。
早期研究集中于内在动机的识别与刻画。
Shah（2006）通过对多个开源社区的质性考察指出，开发者参与的最初驱动力
往往源于对技术问题本身的强烈兴趣——即所谓“抓自己的痒”：
开发者为解决自身实际技术需求而参与，贡献本身即构成收益，
无需外部奖励加以维持\cite{shah2006motivation}。
Davidson 等人（2014）进一步确认，利他主义（即希望为社会留下有价值的成果）
在资深开发者中尤为突出，与兴趣、声誉追求一道构成多层次的内在动机结构\cite{davidson2014older}。

Lerner 和 Tirole（2002）从经济学视角提出了声誉机制假说：
开源社区中代码完全公开，每位开发者的贡献质量对所有人可见，
优秀贡献能够帮助开发者在同行社区中积累声誉，
而该声誉又可转化为劳动力市场上的职业优势\cite{lerner2002simple}。
按此逻辑，声誉实质上是一种“延迟货币”，
表面无偿的贡献可被解读为对人力资本的理性投资。
Hertel 等人（2003）基于对 Linux 内核贡献者的大规模调查，
将开发者动机归纳为实用主义需求、社区认同与政治信念三个维度，
着重论证了集体层面的心理认同对参与行为所具有的独立驱动作用\cite{hertel2003motivation}。
Fang 和 Neufeld（2009）的纵向研究则证实了职业发展动机的重要性——
将开源贡献视为职业发展手段的开发者，其持续参与的稳定性
显著高于纯粹基于兴趣的参与者\cite{fang2009understanding}。
晚近，Tulili 等人（2025）追踪开源生态系统中核心团队的演变过程，
系统刻画了核心开发者流入、流出与留存的动态模式，
确认维持核心团队稳定性的关键在于社区对成员贡献的可见化反馈机制\cite{tulili2025exploring}。

在内在动机与外在动机的理论谱系中，楼天阳等（2014）通过虚拟社区的实验研究报告了一项重要发现：
不同激励政策对成员参与动机的影响差异显著，
物质激励与社会激励的效果取决于社区成员的初始动机结构\cite{lou2014incentive}。
自我决定理论（Self-Determination Theory, Ryan \& Deci, 2000）\cite{ryan2000self}则提供了更为精细的分析框架，
区分了外在动机的不同内化程度：
一旦外在目标被个体充分内化——即个体认为追求该目标与其核心价值观一致——
外在动机便不再与内在动机对立，二者反而协同增强行为动机。
就开源社区而言，若开发者将接受外部资助理解为“对自身技术贡献价值的社会认可”
而非“为金钱而工作”，外在奖励并不必然触发挤出效应；
反之，若奖励被感知为外部控制手段，内在动机遭到侵蚀的可能性则大为增加。
这一分析暗示，捐赠信息的披露方式（究竟以“认可贡献价值”还是“购买服务”
的框架呈现）可能对最终激励效果产生关键影响\cite{benabou2003intrinsic}。

内在动机与外在动机并非简单相加，二者之间的交互关系颇为复杂，
学界以“挤出效应”概括之。
Bénabou 和 Tirole（2003）在理论模型中论证：当个体内在动机足够强时，
引入外在奖励反而可能适得其反——
外在奖励传递了“该工作本身不够有价值”的隐性信号，
同时破坏了个体基于内在动机行动所获得的社会认可\cite{benabou2003intrinsic}。
Qiao 等人（2021）在在线志愿知识贡献平台的实证研究中印证了这一机制：
货币激励短期内确能提升贡献量，
却同时显著降低了高内在动机用户的长期参与意愿\cite{qiao2021mitigating}。
王盼盼等（2022）针对在线健康社区的研究得到了类似结论——
有偿奖励短期内显著提升了医生的知识贡献数量，
但奖励机制撤除后，此前高度依赖内在动机参与的成员
贡献行为出现了显著回落\cite{wang2022reward}。
另一方面，社会性动机同样不可忽视。
Yang 等人（2021）区分了社会影响对初始参与与持续参与的差异化效应，
指出社会认同和群体规范主要驱动开发者作出初始贡献，
持续参与则更多依赖基于互动与信任建立起来的社区归属感\cite{yang2021differential}。
秦敏和李若男（2020）基于在线用户社区的研究同样确认，
社会影响和社区归属感是用户贡献行为形成的关键机制\cite{qin2020contribution}；
秦敏等（2015）更早期的研究则从复杂适应系统理论出发，
阐释了开放式创新社区中用户贡献行为的涌现机制\cite{qin2015cas}。
Sun 等人（2017）报告了货币奖励效果受社会连接性显著调节的证据：
社交关系紧密的平台上，货币奖励与社会认可的结合
能够产生超越单一维度激励的协同效果\cite{sun2017motivation}。
在加密货币生态这一特殊背景下，Kobayakawa 等人（2020）还注意到
代币市场价格表现通过影响开发者对项目前景的预期，
间接左右了开发者的参与强度，表明经济信号在特定情境下
可超越直接货币转移而发挥激励作用\cite{kobayakawa2020impact}。

\subsection{货币激励与捐赠赞助机制的研究综述}

尽管内在动机在开源参与中占据主导地位，随着开源软件战略重要性持续上升，
外部货币激励机制也逐步进入开源生态，引发学界的系统关注。
GitHub 于2019年推出 Sponsors 平台，允许个人和企业直接向开源开发者
支付定期赞助费用，由此构建了迄今最大规模的平台级经济支持基础设施。
Shimada 等人（2022）对该机制开展了最早的系统性研究，报告超过四分之三的
开发者在获得赞助后贡献量有所提升，且部分开发者将接受赞助本身
理解为对开源工作价值的认可，而非单纯的经济交换\cite{shimada2022github}。
Wang 等人（2022b）借助自然实验分析证实，赞助不仅正向影响受赞助开发者的贡献，
还通过社区观察机制对周围开发者产生有限的溢出激励效应\cite{wang2022influence}。
Zhang 等人（2022）考察了“谁赞助谁”的决定因素，确认赞助机制
更倾向于奖励已在社区中建立声誉的活跃开发者，因而更多承担
贡献后的事后认可功能，而非对潜在贡献的事前激励\cite{zhang2022who}。

不过，另一批研究呈现了截然不同的图景。
Overney 等人（2020）对多平台开源捐赠机制的大规模分析指出，
捐赠对开发者活动水平的整体提升作用十分有限——
多数项目即便获得捐赠，开发者的活动模式也并未发生显著变化\cite{overney2020not}。
Conti 等人（2023）追踪 GitHub Sponsors 数年的纵向数据后报告，
实际收到赞助的开发者在长期内创新型贡献有所减少，
经济激励可能将开发者的注意力从开放性技术探索
转向维持赞助关系的稳定性工作\cite{conti2023incentivizing}。
Wang 等人（2022a）在自然实验框架下则观察到显著的
“知识溢出挤出”效应：货币化激励使贡献者减少了免费分享知识的行为\cite{wang2022monetary}。

GitHub Sponsors 研究与非盈利基金会资助研究之间存在一个不容忽视的结构性差异。
GitHub Sponsors 属于“开发者对开发者”或“粉丝对开发者”的点对点直接支持机制，
赞助者自主选择受赞助对象，透明度和即时性均较高；
以 Brink 为代表的专项基金会则充当“中间人”角色，
先集中募资，再经申请审核程序将资金分配给具体开发者。
结构上的差异导致信息传播路径根本不同：
GitHub Sponsors 模式下赞助关系的建立本身即为双方直接互动，
赞助信号直达受赞助者，对旁观者亦具有一定的公开可见性；
基金会模式下，捐款信息的公开传播与资金的最终分配则完全分离——
公告受众远大于最终受益者，该结构天然形成了独立考察信息效应的实验条件。
正是“信息先于货币传播、信息受众大于货币受益者”这一特征，
赋予了基金会模式下信息披露研究超越 GitHub Sponsors 研究的独特价值。
代君等（2019）对开源项目协同开发行为特征的研究亦指出，
开发者之间的协同互动模式对项目成功至关重要，
外部信号对协同行为的激活效应需在结构化的协作网络中加以理解\cite{dai2019collaborative}。

综合以上研究，可识别出一个共同的局限：
现有文献对货币激励的理解普遍聚焦于实际货币转移——即开发者确实收到赞助金额这一事实，
进而将激励效应等同于直接经济补偿的作用。
但捐赠公告所携带的信息价值，
即资金实际到达之前甚至之后，信息本身对未直接受益的开发者所可能产生的激励效应，
在上述研究中均未被单独识别和检验。
捐赠信息披露作为独立于货币转移的激励机制，其有效性究竟如何，
这构成了本研究有别于既有文献的核心问题。

\section{信息披露的理论与实证}

\subsection{自愿信息披露理论}

信息披露研究的系统性发展始于会计学领域，后被管理学和信息系统学科广泛借鉴。
自愿信息披露理论最早由 Verrecchia（1983）
在会计学框架中予以系统阐述\cite{verrecchia1983discretionary}。该理论的核心命题是：
信息不对称环境中，拥有私有信息的组织会权衡披露成本与收益，
自主决定是否以及如何向外部受众传递信息。
一旦披露预期带来的收益（如降低信息使用者的不确定性、
提升对组织的信任、减少信息使用者要求的风险溢价）超过披露成本，
组织便倾向于主动披露\cite{verrecchia1983discretionary}。

对非营利组织而言，信息披露具有维护公信力的特殊意义。
捐赠者和社会公众高度关注非营利组织的资金来源与使用情况，
透明的财务信息披露是此类组织维持捐赠关系、建立社会信任的基础性手段\cite{liu2015nonprofit}。
对 Bitcoin Brink 这类依赖社区捐赠运营的开源资助机构来说，
信息透明度不仅是合规性要求，更是与开发者社区建立信任关系的核心机制。
Brink 在 X 平台公开每笔收到的捐款，在官网公布每位受资助开发者名单——
这种主动且详细的信息披露行为，既是对捐赠者问责的承诺，
也是向开发者社区传递“项目获得外部认可”这一积极信号的渠道。
从自愿披露理论的视角审视，Brink 的披露行为本身即为一种精心设计的信息策略，
其对开发者行为的影响效果值得系统评估。

\subsection{信号理论}

信号理论（Signaling Theory）源自 Spence（1973）
对劳动力市场信息问题所做的经典分析\cite{spence1973job}，
此后被组织行为学、战略管理和信息系统等领域广泛采用。
该理论的基本洞察在于：信息不对称情境下，
处于信息劣势的一方（信号接收者）需依靠可观察的信号
推断自身无法直接观察的信息；
信息优势方（信号发送者）则可通过发送具有可信性的信号，
降低对方的不确定性并引导其作出有利的判断与行动。
有效信号须同时满足两个条件：
其一，与所传递的内容真实相关（相关性）；
其二，具有足够的成本使得虚假信号不值得发送（可信性）\cite{connelly2011signaling}。

在开源捐赠信息披露的情境中，Brink 的两种披露构成了性质迥异的信号。
非定向披露传递的是“项目获得社会认可与外部资金支持”的整体性信号，
类似于众筹平台上项目获得早期资助所产生的“热度信号”——
它向整个开发者群体提供了“该社区正在获得外界关注”的正面判断依据；
定向披露传递的则是“高质量贡献会获得认可与回报”的差异性信号，
通过明确揭示“谁因何贡献而获得资助”强化努力—奖励的因果链条，
向未入选者传递关于资助门槛和评判标准的具体信息，
进而影响其对“自身也有可能获得资助”的预期评估。
两类信号在内容、受众和激活机制上存在本质差异，
这构成了本研究分别检验其激励效应的逻辑起点。

\subsection{奖励信息披露与受众行为}

在信息披露对受众行为影响的实证研究中，众筹领域提供了与本研究最为接近的参照情境。
Mollick（2014）在众筹领域的开创性研究中确认，
项目信息的披露质量与透明度是决定众筹成败的关键因素\cite{mollick2014dynamics}。
Kim 等人（2022）进一步指出，风险信息的披露方式
显著影响投资者的参与决策，信息透明度的提升能有效降低感知风险\cite{kim2022risk}。
Conti 等人（2025）在创业融资情境中论证了另一条路径：
初创企业对开源社区的参与信息能向投资者传递创新能力信号，
信息透明度是连接资金提供方与接收方信任关系的核心机制\cite{conti2025beefing}。
Qiu 和 Zhang（2025）以精细加工可能性模型（ELM）为分析框架，
报告了众筹平台上元信息线索（如项目更新频率、支持者数量等）
通过中心与边缘两条路径显著影响潜在支持者参与决策的证据——
即便这些线索本身并不伴随任何即时经济回报\cite{qiu2025crowdfunding}。
该机制表明，信息本身可作为独立于货币的激励因素发挥作用，
核心路径在于通过改变受众对参与价值的预期评估来引导行为。

知识贡献平台的研究同样提供了重要证据。Toubia 和 Stephen（2013）区分了
“内在效用”（源自贡献内容本身的满足）与“形象效用”
（源自贡献被他人看见和认可的满足），
指出形象效用在很大程度上取决于平台所传递的“贡献受重视”的社会信号\cite{toubia2013intrinsic}。
平台通过公告、排行榜或特殊徽章等方式显式传递此类信号时，
用户参与意愿受到显著影响——尽管这些信号并未带来任何直接金钱收益。
Du（2023）在奖励型众筹场景中的研究进一步证实，
项目信息披露的详细程度直接影响支持者的投资意愿\cite{du2023impact}；
Bai 等人（2024）在金融信息披露领域也注意到，
平台提供的投资者信息透明化机制能显著改善资产定价效率\cite{bai2024platform}。
以上证据共同表明，信息披露能够通过激活受众期望、修正其对奖励可得性的判断，
产生独立且可测量的行为效应，为本研究核心假设提供了实证支撑。

从比较视角看，Brink 的捐赠信息披露与上述研究情境兼有共性与差异。
共性方面，信息的公开传播先于或独立于实际经济转移，
且信息受众（所有关注 Brink 的开发者）远大于实际经济受益者（获得资助的4名开发者）。
差异方面，开源开发者贡献的回报机制更为复杂：
他们面临的不仅是“参与”的二元决策，还有“投入多少”的连续决策，
且该决策嵌入在多年社区参与历史中，受声誉、职业发展和内在动机的交织影响。
这意味着 Brink 信息披露的效应可能比众筹场景中的进展公告效应更为细微、更具情境性。
但分析价值恰在于此：若在如此复杂的动机结构中，
信息披露仍能产生可识别的增量效应，则说明信息因素对开源开发者行为具有
超越噪音水平的独立影响力——这一发现将具有重要理论意涵。

\section{期望理论与期望失验理论}

\subsection{期望理论}

期望理论由心理学家 Victor Vroom 于1964年在工作动机研究中
正式提出\cite{vroom1964work}，是组织行为学和管理学领域引用最为广泛的动机理论之一。
该理论的核心命题在于：个体行为动机并非由需求本身决定，
而取决于其对行动结果的认知预期——具体地说，个体愿意付出努力的程度
由三个认知因素的乘积决定。
其一为期望值（Expectancy），即个体相信自身努力
能达到预期绩效水平的主观概率；
其二为工具性（Instrumentality），即个体相信达到特定绩效
能